\documentclass[10pt,a4paper]{amsart}
\usepackage[margin=19mm]{geometry}
\usepackage{amsmath,amssymb,mathtools}
\usepackage[scr=rsfs,cal=euler]{mathalfa}
\usepackage{xfrac}
\usepackage[unicode,hidelinks,bookmarks=false]{hyperref}
\newcommand{\ie}{i.e.\ }
\newcommand{\cf}{cf.\ }
\newcommand{\resp}{respectively\ }
\newcommand{\Subsection}{\subsection}
\providecommand{\nobreakdash}{\nobreak\hbox{-}}
\setlength{\emergencystretch}{2em}
\multlinegap0pt
\usepackage{bm}
\usepackage{bbm}
\usepackage{dsfont}
\usepackage{xspace}
\usepackage{cancel}
\usepackage{mathtools}
\usepackage[shortlabels]{enumitem}
\usepackage{amsthm}
\makeatletter
\@ifundefined{theorem}{\newtheorem{theorem}{Theorem}}{}
\@ifundefined{claim}{\newtheorem{claim}[theorem]{Claim}}{}
\@ifundefined{lemma}{\newtheorem{lemma}[theorem]{Lemma}}{}
\makeatother
\DeclareMathOperator{\conv}{Conv}
\newcommand{\bbN}{\mathbb{N}}
\newcommand{\bbZ}{\mathbb{Z}}
\newcommand{\bdelta}{{\boldsymbol\delta}}
\newcommand{\calB}{\mathcal{B}}
\newcommand{\calN}{\mathcal{N}}
\newcommand{\calS}{\mathcal{S}}
\newcommand{\ssm}{\smallsetminus}
\title[Average Distortion of Commensurators of Hyperbolic Groups]{Average Distortion of Commensurators of Hyperbolic Groups}
\author{Nir Lazarovich, Suraj Krishna M S, Mahan Mj}
\date{Source: arXiv:2606.27085v1 (CC BY 4.0)}
\begin{document}
\maketitle

\noindent\emph{Source and licence.} Average Distortion of Commensurators of Hyperbolic Groups. Version arXiv:2606.27085v1 (CC BY 4.0), distributed under the Creative Commons Attribution 4.0 International licence, as recorded in the arXiv licence metadata for this exact version and shown on the abstract page https://arxiv.org/abs/2606.27085v1. Full rights notice: licenses/arxiv-2606.27085-CC-BY-4.0.txt.\par\smallskip

\noindent\textit{Editorial scope.} Counted unit: the Lemma whose source label is \texttt{approximating tree} in the article (source0.tex lines 1987--1994, source label \texttt{approximating tree}), delivered together with the article's own complete original proof of it (source0.tex lines 1997--2067, one \texttt{proof} environment of 71 lines). No mathematics is written, repaired or re-derived: statement and proof are the article's own text line for line. Required context retained uncounted: none. Every definition, display and label of those lines is retained unchanged. Omitted from this package and not redistributed: 1--1986, 1995--1996, 2068--2764. Nothing inside a retained excerpt is omitted or altered: every retained body is delivered exactly as the article's lines are written, so no omission receipt of any kind accompanies this package. Citations: the retained excerpts carry no citation command at all, so no bibliography entry of the article is transcribed and no reference is redistributed. Reference adaptation: none was needed and none was made; every reference inside the delivered unit resolves inside it, so no unresolved reference or citation is produced. Printed slips kept literal: no printed word, symbol, hypothesis or number of the counted statement or of its proof was altered. Layout and class: the article's own class is \texttt{amsart} with packages including amsmath, amssymb, amsfonts, xfrac, MnSymbol, amsthm, cite, hyperref, todonotes, enumitem, graphicx, tikz-cd, tikz, color; the wrapper is the lane's pinned \texttt{amsart} editorial wrapper at 10pt on A4, which restates the theorem-like environments the retained text needs and adds the article's own macro definitions that the retained text needs (\texttt{\textbackslash bbN}, \texttt{\textbackslash bbZ}, \texttt{\textbackslash bdelta}, \texttt{\textbackslash calB}, \texttt{\textbackslash calN}, \texttt{\textbackslash calS}, \texttt{\textbackslash conv}, \texttt{\textbackslash ssm}), with extra wrapper packages (bm, bbm, dsfont, xspace, cancel, mathtools, enumitem). The delivered numbering is the unit's own. Assets: no retained span contains a figure environment, a graphics-inclusion command, a PDF-page inclusion, a table or a TikZ picture; no image file is copied into this package, the package declares no asset dependency and no image file is redistributed with it. Licence: version arXiv:2606.27085v1 of this article is distributed under the Creative Commons Attribution 4.0 International licence, as recorded in the arXiv licence metadata for this exact version and shown on the abstract page https://arxiv.org/abs/2606.27085v1. A full rights notice is delivered at licenses/arxiv-2606.27085-CC-BY-4.0.txt. Characters: every retained excerpt is pure ASCII, so the delivered engine, XeLaTeX with its default Latin Modern text font, reports no missing character.
    \begin{lemma}\label{approximating tree}
        Let $A$ be an $M$-barycentrically thick set. Then, there exists a finite tree $T$ and $a_0\in V(T)$ such that:
    \begin{enumerate}[label=(T\arabic*)]
        \item \label{T has no deg 2} All vertices $V(T)\setminus \{a_0\}$ do not have degree 2 -- that is, $V(T) \setminus \{a_0\} \subseteq  V_1(T)\cup V_{\ge 3}(T)$,
        \item \label{leaves of T}  $|V_1(T)|\le |A|$, and
        \item \label{T covers C}  $|\conv(A)| \le \bdelta(M)|V_{\ge 3}(T)|$.
    \end{enumerate}
    \end{lemma}

\begin{proof}
    Let $A$ be an $M$-barycentrically thick set, and let $C = \conv(A)$.
    Fix some $a_0\in A$.
    
    For $R> 0$, let $n:C\ssm\{a_0\} \to \bbZ_{\ge 0}$ be the map defined by $n(x) = \lceil\tfrac{d(x,a_0)}{R}\rceil$ for $x\in C\ssm \{a_0\}$. That is, $n(x)$ is the maximal $n\ge 0$ such that $nR < d(x,a_0)$.

    Consider the set $S_n = \calS(a_0,n\cdot R)\cap C$ for all $n\ge 0$, and their union $S = \bigcup_n S_n$. Let $p:C\ssm\{a_0\}\to S$ be the map that sends $x$ to the unique point of intersection of a geodesic $[a_0,x]$ with $S_{n(x)}$,  i.e.\ $p(x)$ is the nearest point projection of $x$ to a strictly smaller sphere in the collection $S_n$.
    Extend $p$ to a map $p:C\to S$ by setting $p(a_0)=a_0$.

    For $r> 0$, choose a maximal set $V_n$ in $S_n$ such for all $x, y\in V_n$ if $x\ne y$ then $\calB(x,r)\cap \calB(y,r)=\emptyset$.
    Since $V_n$ is maximal, every point in $S_n$ is at distance at most $2r$ from some point in $V_n$.
    Let $V = \bigcup_n V_n$ and let $\pi:C  \to V$ be the map that sends $x\in C\ssm \{a_0\}$ to a point in $V_{n(x)}$ at distance $2r$ away from $p(x)$.

    The following claim shows that iterative applications of $p$ and $\pi$ are uniformly close.

    \begin{claim}\label{claim: distance inequality p pi}
        For $r'\ge r\ge \bdelta$, and $R\ge \bdelta(r')$, if $x,y\in S_n$ and $d(x,y)\le r'$ then 
        \begin{align}
            \label{eq: d(p,p)}d(p(x),p(y)) &\le \bdelta,\\
            \label{eq: d(p,pi)}d(p(x),\pi(y))&\le 2r+\bdelta\\ 
            \label{eq: d(pi,pi)}d(\pi(x),\pi(y))&\le 4r+\bdelta, \text{ and}\\
            \label{eq: d(p j,pi j)} d(\pi^j(x),p^j(x))&\le r',\;  \forall j\in \bbN,
        \end{align}
        Moreover, if $a\in A$, $x\in V_n$,
        \begin{equation}\label{eq: close p implies equal pi}
        \text{ if }d(p^i(a),x')\le r/2 \text{ then } \pi^i(a)=x'
        \end{equation}
        % In particular, for all $i$, and $x\in S_n$, $d(\pi^i(x),p^i(x))\le r'$.
    \end{claim}
    \begin{proof} 
        Indeed, since triangles are thin, and $R\gg r'$ we have that if $d(x,y)\le r'$ then $d(p(x),p(y))\le \bdelta$. If $r' \ge \bdelta(r)$ we also have $$d(p(x),\pi(y))\le 2r+\bdelta \le r'$$
        and also 
        \begin{equation*}
            d(\pi(x),\pi(y))\le 4r+\bdelta\le r'. 
        \end{equation*}
        Iterating these gives \eqref{eq: d(p j,pi j)}.
        
        Finally, to prove \eqref{eq: close p implies equal pi}, let $a\in A$ and assume that $d(p^i(a),x')\le r/2$ for $x'\in V_n$. Then, by \eqref{eq: d(p j,pi j)} $d(p^{i-1}(a),\pi ^{i-1}(a))\le r'$. By \eqref{eq: d(p,p)}, $d(p^i(a),p(\pi^{i-1}(a)))\le \bdelta$. Therefore, for $r\ge \bdelta$, $d(x',p(\pi^{i-1}(a)))\le r/2+\bdelta\le r$. 
        It follows by the definition of $\pi$ that $\pi^i(a)=x'$.
    \end{proof}

    Let $V' = \bigcup_{n\ge 2} \pi^n(A)$, and set $A' := A\ssm\{a_0\}$.
    Consider the tree $T$ whose set of vertices is the abstract disjoint union $A'\sqcup V'$, and edges connect $x \in A'$ to $\pi^2(a)\in V'$ and the vertex $x\in V'$ to $\pi(x)\in V'$.
    
    Clearly, $T$ is a tree.

    We claim that for $r\ge \bdelta(M)$ and $R\ge \bdelta(r)$ the tree $T$ satisfies \ref{T has no deg 2}-\ref{T covers C}.

\bigskip

    \noindent\textbf{\ref{T has no deg 2} and \ref{leaves of T}:}
    Clearly every vertex in $A'$ is a leaf. 
    It remains to show that every vertex of $V' \ssm \{a_0\}$ has degree 3.
    Let $x\in V'\setminus \{a_0\}$, then since $A$ is $M$-barycentrically thick, $x$ is an $M$-barycenter of some $a_1,a_2,a_3\in A$. 
    By considering the $\bdelta$-thin quadrilateral $a_0,a_1,a_2,a_3$, we see that $x$ is an $(M+\bdelta)$-barycenter of $a_0$ and two of the three points $a_1,a_2,a_3$. Without loss of generality, $x$ is an $(M+\bdelta)$-barycenter of $a_0,a_1,a_2$.
    It follows that $d(x,p^{i_1}(a_1))\le M+\bdelta$ and $d(x,p^{i_2}(a_2))\le M+\bdelta$ for some $i_1,i_2\in \bbN$.
    By \eqref{eq: close p implies equal pi}, for appropriate $R\gg r\gg M+\bdelta$ we have $x  = \pi^{i_1} (a_1)=\pi^{i_2}(a_2)$.
    % \nir{this requires some explanation}
    By definition of $V'$ we may assume that at least one of $\max\{i_1,i_2\}\ge 2$.
    For $j=1,2$, let $$y_j = \begin{cases} \pi^{i_j-1}(a_j) & i_j\ge 3\\ a_j &i_j =1, 2. \end{cases}$$ 
    Note that $d(y_1,y_2)>R$, so that $y_1\ne y_2$. 
    This shows that $x$ has degree three, since $y_1,y_2,\pi(x)$ are three distinct neighbors of $x$.

    \bigskip

    \noindent \textbf{\ref{T covers C}:} 
    Let $x\in C$, then $x\in [x_-,x_+]$ for some $x_-,x_+\in A$. Since triangles are thin, $x$ is at distance $\delta$ from a geodesic $[a,a_0]$ where $a\in \{x_-,x_+\}\subseteq A$.
Therefore, $x$ is at distance at most $2R+\delta$ away from some point $p^i(a)$ for $i\ge 2$. By \eqref{eq: d(p,pi)} $d(p^i(a),\pi^i(a)) \le r'$. 
Therefore, $C \subseteq \calN_{R+\delta+r'}(V')$.
    It follows that $|C|\le \bdelta(R)|V'| $.
\end{proof}

\end{document}
